\documentclass[../main.tex]{subfiles}
\graphicspath{{\subfix{../images/}}}

\begin{document}
    \section{Análisis de \texorpdfstring{\emph{L\textsubscript{8}}}{L8}}
    \[
        \boxed{L_8=\{a^{n+1}(aba)^nc^2:n\ge0\}}
    \]

    Este lenguaje crece respecto al exponente \(n\). Si analizamos las cadenas generadas por \(n\in\{0,1,2\}\) obtenemos:
    \begin{align*}
        w_0 &= \text{``acc''}\\
        w_1 &=\text{``aaabacc''}\\
        w_2 &=\text{``aaaabaabacc''}
    \end{align*}

    De esta forma, observamos que cualquier cadena \(w_i\in L_8\) debe iniciar con una ``a'' y terminar con ``cc''.
    Diseñamos la máquina de Turing \(M_8\) para que primero verifique que estas condiciones mínimas se cumplan y luego, a través de reglas de búsqueda, consuma de izquierda a derecha el resto de la cadena. \nolinebreak\(M_8\) llegará a su estado de aceptación si y solo si las subcadenas ``a'' y ``aba'' se repiten \(n\) veces tal como indica \(L_8\).
\end{document}
